\subsection{13.19: a small subdirect-quotient counterexample}
\label{prior:13.19}

The question asks whether \(Q/H\) must be a regular \(p\)-group
when \(H\lhd Q\) and both are subdirect in the same finite
direct product. An earlier construction of Kearnes--Mayr--Ru\v{s}kuc
already implies a negative answer \cite{kearnes-mayr-ruskuc}.
Specifically, their Section~6 has a finite perfect group \(G=HM\),
with \(M\) a 2-group and normal \(N\le M\), satisfying
\([N,N]>[G,N,N]\). In their construction with \(d=2\),
Theorem~4.2 gives \(S=\Gamma_1(G,N)\le G^4\) and
\(T=\Gamma_3(G,N)=\gamma_3(S)\), with \(S/T\) nonabelian
of class two. Both contain the diagonal \(G\), hence are
subdirect. All nondiagonal generators lie in \(M^4\);
since the diagonal is killed in \(S/T\), this quotient is
a finite 2-group. A nonabelian 2-group is irregular.
This is the prior implication used to remove 13.19 from the count.

The experiment's smaller example is transparent.
Let \(D=\operatorname{UT}_3(\F_2)\), writing its elements as
\[
 d(a,b,c)d(a',b',c')
   =d(a+a',b+b',c+c'+ab').
\]
It is dihedral of order eight. In \(D^4\) put
\[
 Q=\{(d(a_i,b_j,c_{ij}))_{1\le i,j\le2}:
                 a,b\in\F_2^2,\ C\in M_2(\F_2)\}.
\]
Its multiplication on the eight coordinates is
\[
 (a,b,C)(a',b',C')=(a+a',b+b',C+C'+a(b')^{\mathsf T}),
\]
so \(|Q|=256\). The map
\[
 \phi(a,b,C)=d(a_1+a_2,b_1+b_2,\textstyle\sum_{i,j}c_{ij})
\]
is onto and multiplicative, since the sum of the entries of
\(a(b')^{\mathsf T}\) is
\((a_1+a_2)(b'_1+b'_2)\).
Its kernel \(H\) has order \(32\) and contains the entire
diagonal \(D\). Both \(H\) and \(Q\) thus project onto all
four factors, with \(Q/H\cong D\).
For \(s=d(1,0,0)\), \(t=d(0,1,0)\), one has
\(s^2=t^2=1\), but \((st)^2=d(0,0,1)\ne1\).
Every derived-subgroup square is trivial, so no correction
allowed in the regularity identity can repair this failure.
The original question permits \(p=2\).

There is also a direct universal construction: every finite
\(p\)-group \(K\) is such a quotient when the factors may vary.
Let \(m=|K|\), choose generators \(g_1,\ldots,g_d\), and let
\(I\) be the augmentation ideal in \(\F_p[K]\).
Induction on \(|K|\) gives \(I^m=0\): choose central \(z\) of
order \(p\); the kernel of the map to
\(\F_p[K/\langle z\rangle]\) is \((z-1)\F_p[K]\), whose
\(p\)th power is zero, and use induction in the quotient.
Let \(A\) be the nonunital free associative algebra on symbols
\(x_{ij}\), \(1\le i\le m,1\le j\le d\), truncated at degree
\(m\). In \(1+A\) take \(Q=\langle1+x_{ij}\rangle\).
Substitution \(x_{ij}\mapsto g_j-1\) maps \(Q\) onto \(K\).
For each \(i\), delete the \(i\)th block of symbols and let
\(G_i\) be the image of \(Q\). These maps jointly embed \(Q\):
each surviving monomial has length less than \(m\), omits some
block, and retains its coefficient when that block is deleted.
The kernel \(H\) of the map onto \(K\) projects onto every
\(G_i\), since
\((1+x_{kj})(1+x_{ij})^{-1}\in H\) projects to \(1+x_{kj}\)
for \(k\ne i\).
This proves the assertion, including simultaneous subdirectness.
The trivial \(K\) is handled by \(Q=H=C_p\).
The construction is not claimed to be small or new.
See \artifact{research/13.19-proof.md},
\artifact{research/13.19-universality.md}, and the source comparison
in \artifact{research/13.19-review.md}.

\subsection{17.34: a bounded-class separation proof using PBW}
\label{prior:17.34}

Let \(\mathcal N_c\) be the full quasivariety of torsion-free
nilpotent groups of class at most \(c\). If \(H\le G\in\mathcal N_c\)
is divisible, then its dominion is exactly \(H\).
The fixed-class strong amalgamation theorem of
d'Elb\'ee--M\"uller--Ramsey--Siniora already implies this
\cite[Theorem~4.35]{delbee-amalgamation}.
Use the lower-central filtration on the rational Lie algebra
of the Malcev completion of \(G\), and its intersection
filtration on the subalgebra corresponding to \(H\).
These are compatible Lazard series even if the latter is
not the intrinsic lower-central series of the subalgebra.
Strong amalgamation separates the two copies outside \(H\).
Finite-dimensional amalgamation suffices in arbitrary rank:
every finite set of constraints lies in a finitely generated,
hence finite-dimensional, nilpotent Lie subalgebra; compactness
then supplies the full pair of embeddings.

Here is the retained alternative proof.
For a rational Lie algebra \(L\) of class at most \(c\) and
subalgebra \(\mathfrak h\), let \(U=U(L)\), \(I\) its
augmentation ideal, and \(J=U\mathfrak h\), a left ideal.
Set
\[
 M=I/(J+I^{c+1}),\qquad d(x)=x+J+I^{c+1}.
\]
This is a left \(U\)-module and
\(d([x,y])=xd(y)-yd(x)\).
Choose a basis of \(L\) adapted both to the lower-central
filtration and to \(\mathfrak h\), weighting a basis element
by its largest filtration index. Order the basis with
the complementary vectors before the \(\mathfrak h\)-vectors.
PBW monomials of weight at least \(n\) span exactly \(I^n\).
One inclusion follows since a vector of weight \(w\) lies
in \(I^w\); the other follows by reordering, because brackets
do not decrease total weight.
Also \(J\) is spanned exactly by ordered monomials containing
a nonempty terminal \(\mathfrak h\)-block: multiplication on
the right by an \(\mathfrak h\)-vector only reorders that block.
Thus the complementary single-letter vectors all survive
independently in \(M\), while the \(\mathfrak h\)-vectors vanish.
This proves \(\ker d=\mathfrak h\), in any dimension.

In \(P=L\ltimes M\), with \(M\) abelian, the two maps
\[
 f(x)=(x,0),\qquad g(x)=(x,d(x))
\]
are injective Lie homomorphisms agreeing exactly on
\(\mathfrak h\). Induction gives
\(\gamma_r(P)\subseteq\gamma_r(L)+I^{r-1}M\).
Since \(I^cM=0\), the class remains at most \(c\).
The numerator \(I\) is essential: replacing it by \(U\)
would generally allow an unwanted extra nilpotency step.

Embed \(G\) in its rational Malcev completion and use the
BCH correspondence, which preserves the class bound and
works in arbitrary rank \cite{bahturin-olshanskii}.
Unique roots make the divisible subgroup \(H\) a rational
Lie subalgebra there. The BCH group of \(P\) lies in
\(\mathcal N_c\), and the restrictions of \(f,g\) to \(G\)
agree exactly on \(H\). This is precisely the required
dominion separation. The result concerns the full
\(\mathcal N_c\), with no assertion for every proper
subquasivariety. The class-two predecessor is Shakhova's
explicit 2015 result \cite{shakhova-2015}.
PBW separation in unrestricted Lie algebras is also prior
work of Bergman \cite{bergman-lie}.
See \artifact{research/17.34-proof.md} and
\artifact{research/17.34-known-consequence.md}.
